% =====================================================================
% paper/sections/cdh_echo_synthesis.tex
%
% PROVENANCE. Berkeley B-SYNTH row 16 (iter148, 2026-07-04):
% CDH Echo on Pillar 4 -- a single-mechanism cross-pillar synthesis.
% Driver: scripts/berkeley/cdh_echo_synthesis.py
% Inputs:  experiments/results/length_bias_iter136_step_coupling.tsv
%          experiments/results/berkeley/cdh_gradnorm_vs_reward.tsv
% Outputs: experiments/results/berkeley/cdh_echo_{pooled_paired,
%          cross_pillar,ratio,sign_test}.tsv
%          experiments/results/berkeley/cdh_echo_summary.json
%
% Cross-pillar synthesis: row 12 (CDH, Pillar-1) is corroborated by
% iter 136 (length-bias, Pillar-4). Same predicted direction holds
% across both pillars (added component LOOSENS |rho| of natural
% reward-vs-axis coupling); the learned component amplifies 3.82x
% more than the fixed component, consistent with CDH's mechanism.
% =====================================================================

\subsection{Cross-Pillar Synthesis: The CDH Echo on Length-Bias Data}
\label{sec:cdh-echo}

The Critic-Degeneracy Hypothesis (introduced in
\ifcsname r@sec:frontier-critic-degeneracy\endcsname Section~\ref{sec:frontier-critic-degeneracy}\else the companion Critic-Degeneracy study\fi) makes a sharp
mechanistic prediction: \emph{any} component added on top of the
stateless group-mean baseline shifts the natural gradient-flow
coupling in the same direction, and the shift is \emph{larger
when the component is learned}. We test this on the Pillar-4
length-bias data (the step-coupling TSV named below) by replacing
the learned value head with Dr.GRPO's \emph{fixed} normalisation
\citep{tong2025drgrpo}: advantages $A_i = r_i - \bar r_g$ with
\emph{no} division by the group reward standard deviation
$\sigma_r$, and a loss whose per-sequence length denominator is
the constant $1/L_{\max}$ rather than $1/|o_i|$.

\paragraph{Data.}
We reuse the iter-136 step-level coupling TSV
(\texttt{experiments/results/length\_bias\_iter136\_step\_coupling.tsv})
with $n_{\text{cell}} = 8$ paired observations across two tasks
(arithmetic\_easy: 5 seeds; gsm8k\_cot: 3 seeds) and two algorithms
(GR / Dr.GR). The relevant variable is
$|\rho(\Delta r_t,\Delta L_t)|$---the per-run reward-velocity /
length-velocity coupling---which is the length-axis analogue of
the $|\rho(\|g_t\|, \mathrm{reward}_t)|$ gradient-reward coupling
tested on Pillar-1 in \ifcsname r@sec:frontier-critic-degeneracy\endcsname Section~\ref{sec:frontier-critic-degeneracy}\else the companion same-stack benchmark\fi.

\paragraph{Pre-registered criterion.}
We test the null of \emph{no shared-sign loosening}: that adding
a component on top of the group-mean baseline does not reduce
$|\rho|$ in both pillars. Across the 8 (task, seed) cells, pooled
$|\rho|$ for Dr.GR vs GR should satisfy (i) the Dr.GR/GR ratio is
below $1.0$---i.e.\ the fixed normaliser \emph{loosens} $|\rho|$,
in the same direction as the learned PPO critic (which loosens it
by $19.5\%$ in Pillar-1), but by less---and (ii) the Pillar-1 and
Pillar-4 effects share sign. The no-shared-sign-loosening null is
rejected (descriptively, at this $n$) if (i) and (ii) both hold.

\paragraph{Result.}
The pooled Dr.GR / GR ratio is $\mathbf{0.949}$ (Dr.GR
$|\rho|=0.358$ vs GR $|\rho|=0.377$, $\Delta=-5.1\%$)---favouring
the fixed-normaliser prediction. Cross-pillar, the Pillar-1
\emph{learned} effect is $-19.5\%$ and the Pillar-4 \emph{fixed}
effect is $-5.1\%$ (negative $\Rightarrow$ added component
loosens $|\rho|$). Both effects share sign; the learned / fixed
effect-size ratio is $\mathbf{3.82\times}$, consistent with the
CDH mechanism's prediction that \emph{learned} components carry
parameter noise beyond what \emph{fixed} components contribute.

\paragraph{Verdict.}
At $n{=}8$ paired cells this is a directional, descriptive
result, not a validated one: both the learned (Pillar-1) and
fixed (Pillar-4) added components loosen their respective
couplings, and the learned effect is the larger of the two
($3.82\times$), as the CDH mechanism predicts---on a
\emph{different pillar, different algorithm pair, and different
coupling axis} (length instead of reward). Pillar-1's PPO/GRPO
and Pillar-4's Dr.GR/GR comparisons are thus \emph{consistent
with} a single mechanism projected onto two coupling axes,
scoped to the small-$n$ evidence available here.
